\BeleskeID{08-s063}
% element: 08-s063:e03
% element: 08-s063:notes
Diferencijator na uzlazni skok daje pozitivan izlaz koji opada prema nuli. Na silazni skok pojavljuje se negativna promena iste amplitude $V_{OH}-V_{OL}$, bez obzira na to koliko je prethodni izlaz stigao da opadne.

Za dugačak impuls prvi odziv skoro iščezne pre drugog skoka, pa se dobijaju dobro razdvojeni pozitivni i negativni odzivi. Za trajanje uporedivo sa $\tau$ drugi skok počinje iz nenulte prethodne vrednosti. Za veoma kratak impuls izlaz se tokom impulsa malo promeni, a posle silaznog skoka ostaje mali negativan rep. Napon otpornika sme da ima skok jer je kontinualan napon kondenzatora, a ne nužno sam izlaz diferencijatora.
